\documentclass[11pt]{article}
\usepackage[margin=1in]{geometry}
\usepackage{amsmath}
\usepackage{amssymb}
\usepackage{graphicx}
\usepackage{siunitx}
\usepackage{hyperref}

\title{Methods --- Millikan Oil Drop Experiment}
\author{}
\date{}

\begin{document}

\section*{Methods}

\subsection*{Experimental Setup}

The Millikan oil drop apparatus (PASCO AP-8210A) was used to measure the electric charge on individual oil droplets suspended between two parallel plates separated by $d = \SI{7.59}{mm}$. The droplets were observed through a telescope with a reticle scale, where each major division corresponded to a distance of \SI{0.500}{mm}. Three video recordings (TRIAL 1--3) of droplet oscillations were captured at approximately \SI{16.68}{fps} and analyzed to extract charge values.

\subsection*{Temperature and Pressure Measurement}

Temperature was measured using a $\SI{10}{k\Omega}$ NTC (negative temperature coefficient) thermistor calibrated against the PASCO apparatus table. The thermistor resistance was recorded as $R = \SI{11200}{-11400}{\Omega}$ across all trials, corresponding to a temperature range of approximately $T = \SI{22.2}{-22.5}{\celsius}$ (nominal $T_{\text{nom}} = \SI{22.35}{\celsius}$). The conversion from resistance to temperature was performed using linear interpolation of the calibration table, with a B-parameter fallback model for out-of-range values:
\begin{equation}
\frac{1}{T} = \frac{1}{T_0} + \frac{1}{B}\ln\left(\frac{R}{R_0}\right)
\end{equation}
where $T_0 = \SI{298.15}{K}$ (25\textdegree C), $R_0 = \SI{10000}{\Omega}$, and $B = \SI{3950}{K}$.

Atmospheric pressure was measured as $P = \SI{75.9}{cmHg} = \SI{101191}{Pa}$ (approximately 1 standard atmosphere). Air density was calculated using the ideal gas law:
\begin{equation}
\rho_{\text{air}} = \frac{PM}{RT}
\end{equation}
where $M = \SI{0.02897}{kg/mol}$ is the molar mass of air and $R = \SI{8.314}{J/(mol \cdot K)}$ is the gas constant. Dynamic viscosity of air was obtained by interpolation of standard tables (PASCO manual Appendix A) at the measured temperature.

\subsection*{Video Analysis and Oscillation Extraction}

High-speed videos were manually analyzed using a custom interactive OpenCV-based frame player. For each observed oil droplet, we recorded the frame numbers and timestamps corresponding to the top and bottom turning points of the droplet's vertical oscillation.

The oscillations occur in two phases:
\begin{itemize}
\item \textbf{Falling phase}: droplet moves downward under gravity and electric force, velocity $v_{\text{fall}}$
\item \textbf{Rising phase}: droplet moves upward (often in electric field), velocity $v_{\text{rise}}$
\end{itemize}

For each phase, the time interval $\Delta t$ between consecutive turning points was computed from video timestamps, and the corresponding distance was calculated as the reticle separation. Thus:
\begin{equation}
v = \frac{d_{\text{reticle}}}{\Delta t}
\end{equation}

Multiple oscillation cycles were recorded per droplet to obtain averaged velocities and error estimates.

\subsection*{Charge Calculation}

The charge on each droplet was calculated using the governing equations from the UCSB Physics 128AL lab manual (PASCO AP-8210A apparatus guide), which account for drag forces and electric forces on the suspended droplet.

\subsubsection*{Radius Estimation}

At terminal velocity falling through air with viscous drag, the droplet radius is estimated from:
\begin{equation}
a_1 = \sqrt{\frac{9 \eta v_{\text{fall}}}{2(\rho_{\text{oil}} - \rho_{\text{air}}) g}} \quad \text{(Eqn.~5)}
\end{equation}
where $\eta$ is the dynamic viscosity of air, $\rho_{\text{oil}} = \SI{853}{kg/m^3}$ is the mineral oil density, $\rho_{\text{air}}$ is the calculated air density, and $g = \SI{9.80665}{m/s^2}$ is gravitational acceleration.

\subsubsection*{Uncorrected Charge}

The charge before Cunningham correction is:
\begin{equation}
q_0 = \frac{18\pi d}{V} \sqrt{\frac{\eta^3 v_{\text{fall}}}{2(\rho_{\text{oil}} - \rho_{\text{air}})g}} (v_{\text{rise}} + v_{\text{fall}}) \quad \text{(Eqn.~6)}
\end{equation}
where $d = \SI{7.59}{mm}$ is the plate separation and $V = \SI{552}{-555}{V}$ is the applied voltage (nominal $V_{\text{nom}} = \SI{553.5}{V}$).

\subsubsection*{Cunningham Correction}

For submicron droplets, the drag coefficient deviates from Stokes law due to gas molecule slip at the droplet surface. This is corrected using the Cunningham correction factor:
\begin{equation}
q = \frac{q_0}{\left[1 + \frac{b}{p a_1}\right]^{3/2}} \quad \text{(Eqn.~10)}
\end{equation}
where $b = \SI{8.226e-3}{Pa \cdot m}$ is the Cunningham constant and $p$ is atmospheric pressure. This correction typically increases the calculated charge by 5--10\% for $a_1 \sim \SI{1}{\mu m}$.

\subsection*{Uncertainty Propagation}

Uncertainties in the measured charge arise from three sources:

\subsubsection*{Systematic Uncertainties (Voltage and Temperature)}

Voltage and temperature were recorded as ranges rather than single values. The charge was evaluated at all four corners of the parameter space $(V_{\text{lo}}, V_{\text{hi}}) \times (T_{\text{lo}}, T_{\text{hi}})$:
\begin{align}
\Delta q_V &= \frac{1}{2}\left[\max_{V \in (V_{\text{lo}}, V_{\text{hi}})} q(V) - \min_{V} q(V)\right] \\
\Delta q_T &= \frac{1}{2}\left[\max_{T \in (T_{\text{lo}}, T_{\text{hi}})} q(T) - \min_T q(T)\right]
\end{align}

The combined systematic uncertainty is:
\begin{equation}
\Delta q_{\text{syst}} = \sqrt{(\Delta q_V)^2 + (\Delta q_T)^2}
\end{equation}

\subsubsection*{Statistical Uncertainty (Velocity Measurements)}

Individual measurements of $v_{\text{fall}}$ and $v_{\text{rise}}$ from the five--seven oscillation cycles per droplet were recorded. Using first-order error propagation and the functional form of Equation~6:
\begin{equation}
\frac{\Delta q}{q} \approx \sqrt{\left(\frac{3}{2}\frac{\sigma_{v_{\text{fall}}}}{v_{\text{fall}}}\right)^2 + \left(\frac{\sigma_{v_{\text{rise}} + v_{\text{fall}}}}{v_{\text{rise}} + v_{\text{fall}}}\right)^2}
\end{equation}

where $\sigma_v$ denotes the standard error of the mean velocity from the sample of oscillations.

\subsubsection*{Total Uncertainty}

Statistical and systematic contributions are combined in quadrature:
\begin{equation}
\Delta q_{\text{total}} = \sqrt{(\Delta q_{\text{stat}})^2 + (\Delta q_{\text{syst}})^2}
\end{equation}

\subsection*{Elementary Charge Estimation}

For each droplet, the measured charge $q$ is a multiple of the elementary charge: $q = n e$, where $n$ is the number of elementary charges. A best-fit integer $n$ was assigned by rounding $|q|/e_{\text{ref}}$ to the nearest integer, where $e_{\text{ref}} = \SI{1.602176634e-19}{C}$ is the CODATA 2018 value.

Droplets within the same charge quantization class ($n$ = constant) were grouped, and the elementary charge was estimated as:
\begin{equation}
e = \frac{\langle q \rangle}{n}
\end{equation}

A weighted average over all droplets, with weights $w_i = 1/\sigma_i^2$ where $\sigma_i$ is the total uncertainty on droplet $i$, was used to obtain a final estimate:
\begin{equation}
e_{\text{est}} = \frac{\sum_i w_i e_i}{\sum_i w_i}, \quad \sigma_e = \frac{1}{\sqrt{\sum_i w_i}}
\end{equation}

\subsection*{Data Output}

For each trial, the following outputs were generated:
\begin{itemize}
\item Per-droplet charge summary CSV file
\item Velocity--time plot for each droplet (showing individual fall/rise velocities with mean lines)
\item Combined analysis CSV (all droplets from the trial)
\item Summary histogram and $q/e$ charge quantization plot
\end{itemize}

\end{document}
